\scanpage{32}
% archive sheet number: 31
% letterhead
\begin{tabular}{c}
INSTITUT\\
DES HAUTES ÉTUDES SCIENTIFIQUES\\
BURES-SUR-YVETTE (S.-\&-O.)\\
928-48-53
\end{tabular}

\uline{Corps de base topologique isom à $\mathbb R$ ou $\mathbb C$}

$K$ corps topologique isom à $\mathbb C$ (il y a \unsure{donc|alors} \uline{deux} isom avec $\mathbb C$, mais on ne choisit pas). $K$ contient can. $\mathbb R$.

$M$ motif sur $K$. Alors \deleted{$\mathrm{DR}(M)$ \unsure{De Rham}} pour \unsure{tt} $X$ lisse projectif sur $K$, $X(K)$ est un espace topologique cpt (variété $K$-analytique \unsure{compact|compacte}), d'où $H^{*}(X,\mathbb Q)$, et on définit ainsi \uline{foncteur de Betti} $M\unsure{\rightsquigarrow|\sim}\mathrm{Betti}(M)$ sur les motifs sur $K$. $\mathrm{Betti}(M)\otimes_{\mathbb Q}K$ \deleted{\unsure{se|$\simeq$}} correspond à $H^*(X,K)$. Par l'argument habituel, ce dernier est isomorphe can. à la coh. de DeR., d'où
\[
 \mathrm{Betti}(M)_K\simeq\mathrm{DR}(M)\qquad \otimes\text{-isom can.}
\]
Ceci met sur \inserted{(le $K$-vecteur)} $\mathrm{DR}(M)$ une structure \uline{réelle}, dont le sous-espace réel est $\mathrm{Betti}(M)_{\mathbb R}$ ; soit $\sigma$ l'involution. La filtration canonique de $\mathrm{DR}(M)$, et $\sigma$, permettent de façon unique de splitter cette filtration, de façon à obtenir sur $\mathrm{DR}(M)$ une \uline{bigraduation} compatible avec le degré total, et redonnant la filtration can. à l'aide du premier degré. Le choix d'un isom $\unsure{\pi|\tau|\sigma|u}:K\unsure{\mathord{-}|\longrightarrow}\mathbb C$ met, \unsure{par} transport de str., une

